\section{Introduction}
\label{sec:intro}

A deep network trained to stage sleep agrees with a human scorer about as well
as two human scorers agree with each other, provided the patient comes from the
cohort it was trained on. Applied to a patient from another clinic, the same
network loses a substantial part of that agreement. The loss is not a failure
of the filters: a slow oscillation is a slow oscillation in every laboratory.
It is a failure of the statistics with which the network represents the
patient, and those statistics are produced, to a large extent, by the
normalization layers.

The default label-free remedy is marginal adaptation. Re-estimating the
batch-normalization statistics on the recording being scored --- adaptive batch
normalization, AdaBN \citep{li2018adabn} --- removes one mean and one variance
per layer and per night, uses no label and changes no weight. It is cheap,
deterministic and widely used. Its premise is that the patient is a
recording-wide offset. This paper shows that the premise is only half true,
and that the half it misses is the one that a single vector per recording
cannot address.

\paragraph{Two parts of the patient}
The class-conditional mean of the representation in a recording decomposes
exactly into three terms (Section~\ref{sec:method}): the stage prototype
$\mu_k$ the classifier was trained on, a recording-wide offset
$a_{\cR}$ that is the same for every epoch of the night, and a
stage-conditional signature $s_{\cR,k}$ that describes how this patient's
wake, N2 or REM differs from the prototype. The post-hoc mean subtraction
that mimics AdaBN at the representation removes $a_{\cR}$ and leaves
$s_{\cR,k}$ untouched; layerwise AdaBN removes $a_{\cR}$ and, as a nonlinear
side effect, in some settings part of $s_{\cR,k}$, but it has no stage index
and cannot target it. The residual is not small either way. Within each stage a linear probe
still identifies the patient far above chance after layerwise AdaBN, and the
conditional part of the residue is what the classifier compares when it
decides between two stages: the logit margin under a marginal correction is
perturbed by $s_{\cR,j}-s_{\cR,k}$.

\paragraph{Stage-conditional shift normalization}
If the shift is indexed by the stage, it can be removed stage by stage. SCSN
estimates, for each stage $k$ and recording $\cR$, the shift
$\hat\delta_{\cR,k}$ between the recording's stage-conditional mean and the
cohort prototype, and subtracts the posterior-weighted average of those shifts
from every epoch:
$\tilde z_i=z_i-\sum_k p_{ik}\hat\delta_{\cR,k}$, where $p_{ik}$ is the
model's own softmax on the uncorrected representation. The estimator is
closed-form with one shrinkage constant, uses no label, changes no weight, and
needs no additional forward pass: the representation and the posterior come
from the same pass over the night. The offset cancels for any
posterior that sums to one; the stage-conditional signature cancels when the
posterior is informative. Under the population correction the class-conditional
means of every recording are mapped back to $\mu_k$, so the logit margin
between any two stages is restored to the one the head was trained on.
Setting the mixing parameter of the interpolated form to zero recovers the
marginal correction, so the two are endpoints of one family rather than
unrelated methods.

\paragraph{Why the distinction matters in practice}
The distinction is not only analytical. A batch-normalized stager can be
adapted by replacing its statistics, and the test-time adaptation literature
has produced a family of increasingly elaborate ways to do so
\citep{li2018adabn,wang2021tent,cafa2023,unmixtns2024,psdnorm2026}. A
group-normalized stager cannot use the standard AdaBN replacement: its
normalization is computed within an epoch and therefore has no test-time batch
statistic to re-estimate. A bespoke entropy objective on GroupNorm's affine
parameters could be constructed, but this is not the usual Tent procedure and
we treat it as a separate method outside this study. Yet group normalization
is an increasingly common choice precisely because it removes dependence on
batch composition, and the representation it produces still carries the
patient. On these checkpoints we compare the uncorrected model, the post-hoc
marginal correction, and SCSN; SCSN is the only stage-aware method among the
statistics-replacement approaches evaluated here.

The same failure separates SCSN from existing transfer methods for sleep and
EEG. Cross-domain sleep staging has been addressed by test-time adaptation
\citep{guo2025tta}, whereas recent cross-subject models use meta-learning,
domain alignment, or self-supervised representation learning
\citep{metasth2026,sedaeeg2025,transformeeg2026}. These approaches improve
subject generalization by changing training objectives or learned
representations. SCSN leaves both the trained filters and the classifier
unchanged and corrects only recording-level statistics that remain after
training.

\paragraph{Measuring the right quantity}
A subject-identification probe fitted to the penultimate representation is
the natural diagnostic, but the overall probe is dominated by the marginal
component: any per-recording centering collapses it, whether or not the
conditional signature survives. We therefore refit the probe \emph{within
each stage}. The within-stage probe is unchanged by a correction that removes
only $a_{\cR}$, and falls only when the correction removes $s_{\cR,k}$. The
distinction is what separates the two families empirically, and it is the
measurement we report throughout.

\paragraph{What kind of contribution this is}
The contribution is analytical before it is algorithmic. The paper adds no
layer, no loss and no training procedure; it characterizes what a trained
stager's representation contains and which inference-time operations can change
which part of it, and the algorithm follows from that characterization. Two
results carry it. The first is an identity: the class-conditional means of a
recording decompose exactly into a cohort prototype, a recording-wide offset
and a stage-conditional signature, and the two parts are identifiable from each
other by construction. The second is an impossibility statement, which we
prove rather than observe: every correction that subtracts one vector per
recording --- an oracle that knows the patient's stage mix, an estimate of the
recording mean, any vector learned from the recording --- leaves the relative
configuration of the stage-conditional means unchanged
(Corollary~\ref{cor:no-marginal}). No better estimate of the marginal statistic
can remove the signature, because the signature is not a statistic of the
recording's mean; only an operation indexed by the stage can reach it. SCSN is
then derived rather than designed: given the posterior, the posterior-weighted
average of the per-stage shifts is the minimum-mean-square-error estimate of
the shift that affects an epoch, and the residual is the posterior mass placed
on the wrong stage (Remark~\ref{rem:mmse}). That it has a closed form is a
consequence of the setting the method is meant for --- a checkpoint whose
weights are frozen and a night that carries no labels --- and the closed form
is the reason the method is deployable, not the reason it is new. Read as a
normalization operator, SCSN is the class-conditional counterpart of AdaBN: its
sufficient statistics are indexed by the classes the model assigns, rather than
by the batch or by the epoch.

\paragraph{What the paper contributes}
\begin{enumerate}
  \item A decomposition of the penultimate representation into a stage
  prototype, a recording offset and a stage-conditional subject signature,
  with propositions that state exactly what marginal adaptation leaves behind
  (Proposition~\ref{prop:marginal}) and what a stage-conditional correction
  removes (Proposition~\ref{prop:conditional}), and a corollary that no
  correction subtracting one vector per recording can alter the relative
  geometry of the stage-conditional means
  (Corollary~\ref{cor:no-marginal}).
  \item SCSN, a closed-form, label-free normalization of the representation
  that needs no retraining and no optimizer, applies unchanged to
  batch-normalized and group-normalized stagers, and contains the marginal
  correction as the endpoint $\lambda=0$ of a one-parameter family
  (Algorithm~\ref{alg:scsn}).
  \item A measurement protocol that separates the two components of patient
  identity, and evidence on two cohorts, two backbones and two normalization
  regimes that layerwise AdaBN leaves a large within-stage signature while
  SCSN removes several times more of it.
  \item Evidence that on group-normalized stagers, where the standard AdaBN
  statistics replacement is undefined, SCSN remains applicable and removes
  the conditional signature; it also improves accuracy and per-patient
  conformal-coverage dispersion over the uncorrected model.
  \item Controls that isolate the mechanism --- mismatched statistics,
  orthogonalised corrections, hard pseudo-labels, low-rank restrictions ---
  a Tent baseline evaluated on the same runs, and a cross-cohort experiment
  that delimits the regime in which a stage-conditional correction is useful.
\end{enumerate}

The method is offline in the same sense as AdaBN: it reads the unlabeled
epochs of the night it is about to score. Clinical staging is performed after
the night, so this is the setting in which a deployed stager operates. No
label, no epoch of another patient and no weight is used.

The paper reports where the method is not the best available one. On
Sleep-EDF SleepCNN, layerwise AdaBN is $1.4$ points more accurate than SCSN,
a difference that is not significant; on Sleep-EDF EEGNet the two are
indistinguishable; and Tent is competitive on the same cohort. On both HMC
backbones SCSN has the highest mean accuracy of the arms, by a margin that is
significant over the uncorrected model in both cases and not significant over
layerwise AdaBN on HMC SleepCNN. The claim defended here is not that
SCSN dominates every baseline on accuracy, but that it removes a component of
patient identity that none of them removes, and that this component can be
measured, predicted and corrected without labels.

\paragraph{Organisation}
Section~\ref{sec:related} places the work among test-time normalization and
subject-variability methods. Section~\ref{sec:method} derives the
decomposition, states the propositions and defines the algorithm.
Section~\ref{sec:setup} describes the cohorts, arms, metrics and statistics.
Section~\ref{sec:results} reports the main comparison, the mechanism analyses,
the ablations, the Tent baseline, the group-normalized setting and
cross-cohort transfer. Section~\ref{sec:discussion} discusses what the results
do and do not license.
